\begin{frame}{확인 문제 (1)}
  \begin{itemize}
    \item \textbf{1.} ``5\%만 양성인 데이터에 신경망을 써서 정확도
      95\%를 냈다''는 발표. 가장 약한 부분은? Week 2.3 개념으로
      반박 가능한 질문 하나를 쓰시오.
      \vskip0.3em
      {\footnotesize \textbf{답}: 정확도의 함정(Ch02.3) ---
      \textbf{모두 음성}으로만 찍어도 95\%. 질문: ``Precision과
      Recall(또는 PR-AUC)이 각각 얼마인가요? 무조건 음성 모델과
      같은 수인데, 실제로 양성을 얼마나 더 잡는지 보여주세요.''}
    \item \textbf{2.} ``\textbf{테스트셋}으로 여러 $\lambda$ 중
      가장 좋은 것을 골랐다''고 인정한 팀. Week 6.3의 어떤
      원칙이 깨졌고, 왜 그 숫자가 낙관적인가?
      \vskip0.3em
      {\footnotesize \textbf{답}: ``test를 하이퍼파라미터 선택에
      쓰지 말고 딱 한 번만'' 원칙 위반. 후보 중 \textbf{가장 잘
      나온 것}을 고르면 잡음의 max를 고르는 것 --- 후보가 많을수록
      낙관적으로 읽힘(선택 편향). 올바르게: val로 $\lambda$
      선택 후 test \textbf{한 번}.}
  \end{itemize}
\end{frame}
